\documentclass[11pt]{article}
\usepackage{amsmath,amssymb}
\begin{document}
\section*{Section 3, Lemma 3.20: deriving the rook matching coordinates}
\textbf{Manuscript frozen.} This separate note proposes only
a local correction to the \(R_3\) cross-degree-one case.

Let \(G=R_3\) be the nine-vertex rook graph. Suppose an
eighteen-vertex EPPA host \(H\) contains the designated
induced copy of \(G\), and a vertex \(w\) outside it
has exactly one neighbour in \(G\).
The nine translations of the \(3\times3\) grid are
genuine automorphisms of \(R_3\). Their orbit on
a singleton neighbourhood has precisely nine elements.
By EPPA all these profiles occur outside the selected
copy. Since there are only nine outside vertices,
each singleton occurs exactly once and no other profile
occurs. The cross graph is therefore a \emph{matching},
not merely a bipartite relation on two halves
which were previously labelled by hand.

Under the resulting canonical equivalence
\(f:V(R_3)\simeq V(H)\setminus V(G)\), every
automorphism of the rook graph has an EPPA lift
which sends \(f(x)\) to \(f(p(x))\).
This needs no coherent family of lifts:
profile uniqueness forces the action.
Since the full rook automorphism group has exactly
two orbits on unordered distinct vertex pairs,
the induced graph on the outside is
\(\varnothing,R_3,\overline{R_3}\), or \(K_9\).
Vertex-transitivity of \(H\) forces the two halves
to have equal numbers of edges (18), leaving only
\(R_3\) and \(\overline{R_3}\).

\paragraph{Lean target.}
\texttt{RookThreeSingletonOutside.lean} uses the
reusable generic EPPA outside-orbit saturation lemma
and explicit \(72\) rook-automorphism action to derive
the actual outside equivalence and two graph forms.
The principal statement is
\texttt{rook\_eppa\_singleton\_outside\_graph\_two\_forms}.
The theorem assumes only one singleton outside profile,
not uniform cross-degree, an outside partition or a
coherent lift action.

The remaining graph-identification step is to construct
the matching-layer equivalence of the whole actual host
and connect its two canonical variants to the named
alternatives in Lemma~3.20.
The original source is unchanged until author review.
\end{document}
